The dynamics of M31 (Andromeda) and the Milky Way (MW) can be used to estimate
the total mass of the Local Group (LG). The basic idea is that galaxies
currently in a binary system were at approximately the same point in space
shortly after the Big Bang. To a reasonable approximation, the mass of the
local group is dominated by the masses of the MW and M31. Via Doppler shifts
of the spectral lines, it was found that M31 is moving towards the MW with a
speed of 118\,km\,s$^{-1}$. This may be surprising, given that most galaxies
are moving away from each other with the general Hubble flow. The fact that
M31 is moving towards the MW is presumably because their mutual gravitational
attraction has eventually reversed their initial velocities. In principle, if
the pair of galaxies is well-represented by isolated point masses, their total
mass may be determined by measuring their separation, relative velocity and
the time since the universe began. Kahn and Woltjer (1959) used this argument
to estimate the mass in the LG.

In this problem we will follow this argument through our calculation as
follows.

\begin{description}
\item[(a)] Consider an isolated system with negligible angular momentum of two
  gravitating point masses $m_1$ and $m_2$ (as observed by an inertial
  observer at the centre of mass).

  \ioaafig{2017_TH_T11-f1.pdf}

  Write down the expression of the total mechanical energy ($E$) of this
  system in mathematical form connecting $m_1$, $m_2$, $r_1$, $r_2$, $v_1$,
  $v_2$, and the universal gravitational constant $G$, where $v_1$ and $v_2$
  are the radial velocities of $m_1$ and $m_2$, respectively. \hfill[5]
\item[(b)] Re-write the equation in a) in terms of $r$, $v$, $\mu$, $M$, and
  $G$, where $r \equiv r_1 + r_2$ is the separation distance between $m_1$ and
  $m_2$, $v$ is the changing rate of the separation distance,
  $\mu \equiv \dfrac{m_1 m_2}{m_1 + m_2}$ is the reduced mass of the system,
  and $M \equiv m_1 + m_2$ is the total mass of the system. \hfill[10]
\item[(c)] Show that the equation in b) yields
  \[ v^2 = (2GM)\left(\frac{1}{r} - \frac{1}{r_0}\right),
     \quad\text{where } r_0 \text{ is a new constant.} \]
  Find $r_0$ in terms of $\mu$, $M$, $G$ and $E$. \hfill[5]
\end{description}

The solution of the equation in b) is given below in parametric form, under
the initial condition $r = 0$ at $t = 0$:
\[ r(\theta) = \frac{r_0}{2}\left(1 - \cos\theta\right), \]
\[ t(\theta) = \left(\frac{r_0^3}{8GM}\right)^{\frac{1}{2}}
   \left(\theta - \sin\theta\right), \]
where $\theta$ is in radians.

\begin{description}
\item[(d)] From the above parametric equations, show that an expression for
  $\dfrac{vt}{r}$ is
  \[ \frac{vt}{r} = \frac{(\sin\theta)(\theta - \sin\theta)}
     {(1 - \cos\theta)^2} \]
  \hfill[10]
\item[(e)] Now we consider $m_1$ and $m_2$ as the MW and M31, respectively,
  such that the current values of $v$ and $r$ are $v = -118$\,km\,s$^{-1}$ and
  $r = 710$\,kpc, and $t$ may be taken to be the age of the Universe (13700
  million years). Find $\theta$ using numerical iteration. \hfill[10]
\item[(f)] Use the value of $\theta$ from e) to calculate the maximum distance
  between M31 and the MW, $r_{\max}$, and hence also obtain the value of $M$
  in solar masses. \hfill[10]
\end{description}
